\documentclass[12pt,a4paper,fontset=none]{ctexart}
\usepackage[a4paper,margin=1.7cm]{geometry}
\usepackage{fontspec,amsmath,array,booktabs,tabularx}
\usepackage{enumitem,fancyhdr,float,hyperref}
\usepackage[most]{tcolorbox}
\usepackage{xcolor}

\hypersetup{colorlinks=true,linkcolor=black,urlcolor=blue}
\setmainfont{Times New Roman}\setsansfont{Arial}\setmonofont{Menlo}
\setCJKmainfont{Songti SC}\setCJKsansfont{Heiti SC}\setCJKmonofont{Songti SC}
\setlength{\emergencystretch}{1.5em}
\sloppy
\hfuzz=15pt
\hbadness=10000
\pagestyle{fancy}\fancyhf{}
\fancyhead[L]{ROSE 实验四}
\fancyhead[R]{Oberon-0 翻译模式设计}
\fancyfoot[C]{\thepage}
\setlength{\headheight}{15pt}
\setlist[itemize]{leftmargin=1.8em,itemsep=0.25em}
\renewcommand{\arraystretch}{1.22}
\newcolumntype{P}[1]{>{\raggedright\arraybackslash}p{#1}}
\newcolumntype{Y}{>{\raggedright\arraybackslash}X}

\definecolor{accent}{HTML}{9A3412}
\definecolor{accenttwo}{HTML}{0F4C81}
\definecolor{accentthree}{HTML}{2F6B3D}
\definecolor{softb}{HTML}{E8F0F7}
\definecolor{linegray}{HTML}{D7CEC5}
\definecolor{boxbg}{HTML}{FCFAF7}

\newtcolorbox{plainBox}{enhanced,colback=boxbg,colframe=linegray,boxrule=0.7pt,arc=1.2mm,left=2.3mm,right=2.3mm,top=1.8mm,bottom=1.8mm,before upper=\raggedright}

\title{Oberon-0 语言的翻译模式设计}
\author{梁力航\quad 23336128}
\date{2026年5月20日}

\begin{document}
\hfuzz=20pt
\hbadness=10000
\maketitle

\begin{plainBox}
本文档对应 ROSE 项目实验四的提交要求：描述 Oberon-0 语言的翻译模式（Translation Scheme），包括文法转换、属性传递方案和递归子程序映射，并比较自顶向下与自底向上两种分析策略。
\end{plainBox}

\section{文法转换}

为适应递归下降预测分析（自顶向下、LL），对原始 Oberon-0 EBNF 文法做了以下转换。

\subsection{消除左递归}

原始 grammar 中三处存在直接左递归，转换为带迭代循环的等价形式：

\begin{table}[H]\centering
\begin{tabularx}{\textwidth}{>{\ttfamily\small}Y>{\ttfamily\small}Y}
\toprule
\textbf{原始（左递归）} & \textbf{转换后（{ }* 迭代）} \\
\midrule
simple\_expr → term \newline
\quad\ \ |\ simple\_expr + term \newline
\quad\ \ |\ simple\_expr - term \newline
\quad\ \ |\ simple\_expr OR term & simple\_expr → term \{ + term | - term | OR term \}* \\
\midrule
term → factor \newline
\quad |\ term * factor \newline
\quad |\ term DIV factor \newline
\quad |\ term MOD factor \newline
\quad |\ term \& factor & term → factor \{ * factor | DIV factor | MOD factor | \& factor \}* \\
\midrule
declarations → ... \newline
\quad |\ declarations proc\_decl ";" & declarations → ... \{ proc\_decl ";" \}* \newline
\quad\quad（用 while 循环实现） \\
\bottomrule
\end{tabularx}
\caption{左递归消除规则}
\end{table}

\subsection{提取左公因子}

Oberon-0 文法已天然避免了大部分左公因子问题。唯一的潜在歧义是 statement 规则中 \texttt{IDENTIFIER} 既可能开始一个赋值语句，也可能开始一个过程调用。通过一个 token 的向前看（lookahead）区分：

\begin{itemize}
  \item \texttt{IDENTIFIER selectors ":=" expr} → 赋值语句
  \item \texttt{IDENTIFIER selectors "(" params ")"} → 过程调用
  \item \texttt{IDENTIFIER selectors}（后无 \texttt{:=} 或 \texttt{(}） → 无参过程调用
\end{itemize}

\section{翻译模式设计}

翻译模式的核心思路：每个非终结符对应一个递归子程序，继承属性用形式参数传递，综合属性用返回值表示。

\subsection{属性传递方案}

\begin{table}[H]\centering
\footnotesize
\setlength{\tabcolsep}{3pt}
\begin{tabularx}{\textwidth}{>{\bfseries}P{3cm}>{\bfseries}P{2.3cm}>{\bfseries}P{2cm}>{\raggedright\arraybackslash}X}
\toprule
非终结符 & 继承属性 & 综合属性 & 语义动作 \\
\midrule
module & — & Module 对象 & 创建 Module，依次添加 Procedure \\
procedure\_decl & Module & — & 创建 Procedure，注册到 Module \\
statement\_seq & Proc./StmtSeq & — & 按序解析 statements \\
if\_statement & — & IfStatement 对象 & 创建 IfStmt，递归填充分支 \\
while\_statement & — & WhileStmt 对象 & 创建 WhileStmt，递归填充循环体 \\
expression & — & String（文本） & 重建表达式文本用于流程图显示 \\
\bottomrule
\end{tabularx}
\caption{核心非终结符的属性传递方案}
\end{table}

\subsection{递归子程序映射}

\begin{table}[H]\centering
\footnotesize
\setlength{\tabcolsep}{3pt}
\begin{tabularx}{\textwidth}{>{\bfseries}P{3.5cm}>{\raggedright\arraybackslash}X}
\toprule
文法非终结符 & 递归子程序方法 \\
\midrule
module & \texttt{parseModule()} \\
declarations & \texttt{parseDeclarations()} \\
procedure\_declaration & \texttt{parseProcedureDeclaration()} \\
statement\_sequence & \texttt{parseStatementSequence(Procedure)} \\
statement & \texttt{parseStatementCommon(Object\ container)} \\
if\_statement & \texttt{parseIfStatement(Procedure)} \\
while\_statement & \texttt{parseWhileStatement(Procedure)} \\
expression & \texttt{parseExpressionText()\ →\ String} \\
simple\_expression & \texttt{while} 循环调用 \texttt{parseTermText()} \\
term & \texttt{while} 循环调用 \texttt{parseFactorText()} \\
factor & \texttt{parseFactorText()\ →\ String} \\
type & \texttt{parseType()\ →\ String} \\
\bottomrule
\end{tabularx}
\caption{完整递归子程序映射（12 个方法）}
\end{table}

\subsection{流程图生成嵌入}

语义动作在解析过程中直接调用 ROSE Flowchart API：

{\raggedright
\begin{itemize}
  \item 赋值/过程调用 → \texttt{add(new\ PrimitiveStatement(text))}
  \item IF 语句 → \texttt{IfStatement(condition)} → 递归填充 \texttt{getTrueBody()\ /\ getFalseBody()}
  \item WHILE 语句 → \texttt{WhileStatement(condition)} → 递归填充 \texttt{getLoopBody()}
  \item ELSIF 链 → 嵌套 IfStatement 到前一个 IF 的 false 分支
\end{itemize}
}

\section{自顶向下 vs 自底向上——分析策略比较}

通过实验三（JavaCUP / LALR）和实验四（递归下降）的亲身实践，对如下维度的体会：

\begin{table}[H]\centering
\begin{tabularx}{\textwidth}{>{\bfseries}P{2.8cm}YY}
\toprule
维度 & 递归下降（本实验） & LALR（实验三） \\
\midrule
分析方向 & 自顶向下 & 自底向上 \\
简单性 & 直观，方法直接映射非终结符，调试时打断点即可跟踪 & 需理解分析表、状态机和移进/归约逻辑 \\
通用性 & 需手工消除左递归，不能处理所有 CFG & 可处理 LR(1) 文法，更通用 \\
语义动作 & 可在解析过程中任意位置嵌入 & 仅在产生式规约时执行 \\
错误恢复 & 较易实现，方法级 panic-mode & 需特殊 error 产生式或分析表条目 \\
分析速度 & O(n)，无表查找开销 & O(n)，有状态表查找开销 \\
表格大小 & 无需分析表 & LR 分析表可能较大 \\
\bottomrule
\end{tabularx}
\caption{六维对比}
\end{table}

具体体会：
\begin{enumerate}
  \item \textbf{流程图生成的语义动作在递归下降中更自然}——解析到 IF 关键字后可立即创建 IfStatement 对象，然后递归解析 true-branch 时直接向其中添加语句
  \item \textbf{左递归消除是递归下降的额外成本}——但 Oberon-0 仅需处理表达式和声明列表 3 处，成本可控
  \item \textbf{错误恢复在递归下降中更灵活}——可直接跳过当前输入直到同步符号（END、分号），不需额外设计
  \item \textbf{LALR 的表驱动方式在文法演进时更鲁棒}——新增语法规则不一定需要修改现有规则的处理逻辑
\end{enumerate}

\end{document}
